% 原书目录：印刷xxi--xxiii／PDF物理22--24.
\clearpage
\section*{原书目录（页码沿用原书）}\phantomsection\label{sec:source-contents}
\addcontentsline{toc}{section}{原书目录（页码沿用原书）}
以下保留原书目录的全部层级、编号与原印刷页码. 标题链接到本译稿的对应内容；右栏页码不表示本译稿页码.
\begingroup
\small
\renewcommand{\arraystretch}{1.08}
\begin{longtable}{@{}p{1.4cm}p{11.2cm}r@{}}
\textbf{编号}&\textbf{标题}&\textbf{原页}\\\hline
\endfirsthead
\textbf{编号}&\textbf{标题（续）}&\textbf{原页}\\\hline
\endhead
&\hyperref[sec:preface]{\textbf{前言}}&vii\\
&\hyperref[sec:introduction]{\hspace{1em}最优运输导论}&xiv\\
&\hyperref[sec:acknowledgments]{\textbf{致谢}}&xix\\
&\hyperref[sec:notation]{\textbf{记号}}&xxv\\
\textbf{1}&\hyperref[sec:1]{\textbf{原问题与对偶问题}}&1\\
1.1&\hyperref[sec:1.1]{Kantorovich 问题与 Monge 问题}&1\\
1.2&\hyperref[sec:1.2]{对偶性}&9\\
1.3&\hyperref[sec:1.3]{$h$ 严格凸时 $c(x,y)=h(x-y)$ 的情形及最优 $\mathrm T$ 的存在性}&13\\
1.3.1&\hyperref[sec:1.3.1]{\hspace{1em}$\mathbb R^d$ 中的二次情形}&16\\
1.3.2&\hyperref[sec:1.3.2]{\hspace{1em}平坦环面上的二次情形}&18\\
1.4&\hyperref[sec:1.4]{存在性的反例}&20\\
1.5&\hyperref[sec:1.5]{Kantorovich 问题作为 Monge 问题的松弛}&21\\
1.6&\hyperref[sec:1.6]{$c$-凹性、对偶性与最优性}&25\\
1.6.1&\hyperref[sec:1.6.1]{\hspace{1em}凸函数与 $c$-凹函数}&25\\
1.6.2&\hyperref[sec:1.6.2]{\hspace{1em}$c$-循环单调性与对偶性}&28\\
1.6.3&\hyperref[sec:1.6.3]{\hspace{1em}对偶性的一个直接证明}&35\\
1.6.4&\hyperref[sec:1.6.4]{\hspace{1em}最优性的充分条件与稳定性}&37\\
1.7&\hyperref[sec:1.7]{讨论}&41\\
1.7.1&\hyperref[sec:1.7.1]{\hspace{1em}概率解释}&41\\
1.7.2&\hyperref[sec:1.7.2]{\hspace{1em}极分解}&42\\
1.7.3&\hyperref[sec:1.7.3]{\hspace{1em}匹配问题与经济学解释}&44\\
1.7.4&\hyperref[sec:1.7.4]{\hspace{1em}多边缘运输问题}&48\\
1.7.5&\hyperref[sec:1.7.5]{\hspace{1em}鞅最优运输与金融应用}&51\\
1.7.6&\hyperref[sec:1.7.6]{\hspace{1em}Monge--Amp\`ere 方程与正则性}&54\\
\textbf{2}&\hyperref[sec:2]{\textbf{一维问题}}&59\\
2.1&\hyperref[sec:2.1]{一维中的单调运输映射与方案}&59\\
2.2&\hyperref[sec:2.2]{单调映射的最优性}&63\\
2.3&\hyperref[sec:2.3]{Knothe 运输}&67\\
2.4&\hyperref[sec:2.4]{作为 Brenier 映射极限的 Knothe 映射}&72\\
2.5&\hyperref[sec:2.5]{讨论}&79\\
2.5.1&\hyperref[sec:2.5.1]{\hspace{1em}直方图均衡化}&79\\
2.5.2&\hyperref[sec:2.5.2]{\hspace{1em}由一维构造得到的单调映射}&81\\
2.5.3&\hyperref[sec:2.5.3]{\hspace{1em}通过 Knothe 或 Brenier 映射证明等周不等式}&83\\
\textbf{3}&\hyperref[sec:3]{\textbf{$L^1$ 与 $L^\infty$ 理论}}&87\\
3.1&\hyperref[sec:3.1]{成本为 $|x-y|$ 的 Monge 情形}&87\\
3.1.1&\hyperref[sec:3.1.1]{\hspace{1em}距离成本的对偶性}&88\\
3.1.2&\hyperref[sec:3.1.2]{\hspace{1em}次级变分问题}&89\\
3.1.3&\hyperref[sec:3.1.3]{\hspace{1em}运输射线的几何性质}&91\\
3.1.4&\hyperref[sec:3.1.4]{\hspace{1em}最优运输映射的存在与不存在}&99\\
3.1.5&\hyperref[sec:3.1.5]{\hspace{1em}单调运输的逼近}&101\\
3.2&\hyperref[sec:3.2]{上确界情形：$L^\infty$}&104\\
3.3&\hyperref[sec:3.3]{讨论}&109\\
3.3.1&\hyperref[sec:3.3.1]{\hspace{1em}不同范数与更一般的凸成本}&109\\
3.3.2&\hyperref[sec:3.3.2]{\hspace{1em}凹成本（$L^p$，$0<p<1$）}&113\\
\textbf{4}&\hyperref[sec:4]{\textbf{最小流}}&121\\
4.1&\hyperref[sec:4.1]{Euler 观点与 Lagrange 观点}&121\\
4.1.1&\hyperref[sec:4.1.1]{\hspace{1em}静态与动态模型}&121\\
4.1.2&\hyperref[sec:4.1.2]{\hspace{1em}连续性方程}&123\\
4.2&\hyperref[sec:4.2]{Beckmann 问题}&127\\
4.2.1&\hyperref[sec:4.2.1]{\hspace{1em}引言、形式等价与变体}&127\\
4.2.2&\hyperref[sec:4.2.2]{\hspace{1em}构造 Beckmann 问题的极小解}&129\\
4.2.3&\hyperref[sec:4.2.3]{\hspace{1em}曲线上的测度的交通强度与交通流}&134\\
4.2.4&\hyperref[sec:4.2.4]{\hspace{1em}一维中的 Beckman 问题}&140\\
4.2.5&\hyperref[sec:4.2.5]{\hspace{1em}最优 $\mathbf w$ 的刻画与唯一性}&142\\
4.3&\hyperref[sec:4.3]{运输密度的可积性}&144\\
4.4&\hyperref[sec:4.4]{讨论}&151\\
4.4.1&\hyperref[sec:4.4.1]{\hspace{1em}拥挤运输}&151\\
4.4.2&\hyperref[sec:4.4.2]{\hspace{1em}分支运输}&162\\
\textbf{5}&\hyperref[sec:5]{\textbf{Wasserstein 空间}}&177\\
5.1&\hyperref[sec:5.1]{定义与三角不等式}&179\\
5.2&\hyperref[sec:5.2]{$W_p$ 诱导的拓扑}&183\\
5.3&\hyperref[sec:5.3]{$\mathbb W_p$ 中的曲线与连续性方程}&187\\
5.3.1&\hyperref[sec:5.3.1]{\hspace{1em}Benamou--Brenier 泛函 $\mathscr B_p$}&189\\
5.3.2&\hyperref[sec:5.3.2]{\hspace{1em}绝对连续曲线具有速度场}&192\\
5.3.3&\hyperref[sec:5.3.3]{\hspace{1em}连续性方程的正则化}&194\\
5.3.4&\hyperref[sec:5.3.4]{\hspace{1em}速度场给出 Lipschitz 性质}&197\\
5.3.5&\hyperref[sec:5.3.5]{\hspace{1em}$W_p^p$ 沿测度曲线的导数}&198\\
5.4&\hyperref[sec:5.4]{$\mathbb W_p$ 中的恒速测地线}&202\\
5.5&\hyperref[sec:5.5]{讨论}&207\\
5.5.1&\hyperref[sec:5.5.1]{\hspace{1em}$W_\infty$ 距离}&207\\
5.5.2&\hyperref[sec:5.5.2]{\hspace{1em}Wasserstein 距离与负阶 Sobolev 距离}&209\\
5.5.3&\hyperref[sec:5.5.3]{\hspace{1em}Wasserstein 距离与分支运输距离}&211\\
5.5.4&\hyperref[sec:5.5.4]{\hspace{1em}切片 Wasserstein 距离}&214\\
5.5.5&\hyperref[sec:5.5.5]{\hspace{1em}$\mathbb W_2$ 中的重心}&215\\
\textbf{6}&\hyperref[sec:6]{\textbf{数值方法}}&219\\
6.1&\hyperref[sec:6.1]{Benamou--Brenier}&220\\
6.2&\hyperref[sec:6.2]{Angenent--Hacker--Tannenbaum}&225\\
6.3&\hyperref[sec:6.3]{Monge--Amp\`ere 方程的数值解}&232\\
6.4&\hyperref[sec:6.4]{讨论}&235\\
6.4.1&\hyperref[sec:6.4.1]{\hspace{1em}离散数值方法}&235\\
6.4.2&\hyperref[sec:6.4.2]{\hspace{1em}半离散数值方法}&242\\
\textbf{7}&\hyperref[sec:7]{\textbf{概率测度上的泛函}}&249\\
7.1&\hyperref[sec:7.1]{半连续性}&250\\
7.1.1&\hyperref[sec:7.1.1]{\hspace{1em}势能、相互作用、运输成本与对偶范数}&250\\
7.1.2&\hyperref[sec:7.1.2]{\hspace{1em}局部泛函}&254\\
7.2&\hyperref[sec:7.2]{凸性、第一变分与次微分}&260\\
7.2.1&\hyperref[sec:7.2.1]{\hspace{1em}对偶范数}&263\\
7.2.2&\hyperref[sec:7.2.2]{\hspace{1em}运输成本}&264\\
7.2.3&\hyperref[sec:7.2.3]{\hspace{1em}最优性条件}&267\\
7.3&\hyperref[sec:7.3]{位移凸性}&269\\
7.3.1&\hyperref[sec:7.3.1]{\hspace{1em}$\mathscr V$ 与 $\mathscr W$ 的位移凸性}&270\\
7.3.2&\hyperref[sec:7.3.2]{\hspace{1em}$\mathscr F$ 的位移凸性}&271\\
7.3.3&\hyperref[sec:7.3.3]{\hspace{1em}广义测地线上的凸性}&275\\
7.4&\hyperref[sec:7.4]{讨论}&276\\
7.4.1&\hyperref[sec:7.4.1]{\hspace{1em}一个案例研究：$\min F(\varrho)+W_2^2(\varrho,\nu)$}&276\\
7.4.2&\hyperref[sec:7.4.2]{\hspace{1em}Brunn--Minkowski 不等式}&279\\
7.4.3&\hyperref[sec:7.4.3]{\hspace{1em}城市均衡}&280\\
\textbf{8}&\hyperref[sec:8]{\textbf{梯度流}}&285\\
8.1&\hyperref[sec:8.1]{$\mathbb R^d$ 及度量空间中的梯度流}&285\\
8.2&\hyperref[sec:8.2]{$\mathbb W_2$ 中的梯度流与偏微分方程的推导}&290\\
8.3&\hyperref[sec:8.3]{Fokker--Planck 情形的分析}&293\\
8.4&\hyperref[sec:8.4]{讨论}&301\\
8.4.1&\hyperref[sec:8.4.1]{\hspace{1em}EVI、唯一性与测地凸性}&301\\
8.4.2&\hyperref[sec:8.4.2]{\hspace{1em}其他梯度流偏微分方程}&304\\
8.4.3&\hyperref[sec:8.4.3]{\hspace{1em}Dirichlet 边界条件}&311\\
8.4.4&\hyperref[sec:8.4.4]{\hspace{1em}演化偏微分方程：不止梯度流}&313\\
&\hyperref[sec:exercises]{\textbf{习题}}&325\\
&\hyperref[sec:references]{\textbf{参考文献}}&339\\
&\hyperref[sec:index]{\textbf{索引}}&351\\
\end{longtable}
\endgroup
